%% kapa: exam=1 round=23 year=2012 period=2 form=A subject=appraisal_law subject_ko=부동산관계법규 seq=15 number=55
%% source: Q-Net 감정평가사 기출문제 (cst00309) artlSeq=5052580 file=raw/1008/5052580/2012년 제23회 감정평가사 제1차시험 기출문제(2교시 A형).pdf
%% conversion: kordoc 4.24.51 -> tools/gen1.py (auto); review=proofread-sonnet; flags=none
\begin{question}{55}
국토의 계획 및 이용에 관한 법령상 도시계획 등에 관한 설명으로 옳은 것은?
\begin{choices}
  \item 도시계획은 광역도시계획과 도시기본계획으로 구분한다.
  \item 개발제한구역, 시가화조정구역의 지정 또는 변경에 관한 계획은 도시기본계획으로 결정한다.
  \item 도시관리계획을 시행하기 위한 정비사업은 도시계획시설사업에 포함된다.
  \item 도시계획시설은 기반시설 중 도시기본계획으로 결정된 시설을 말한다.
  \item 자동차 및 건설기계검사시설과 그 시설의 기능발휘와 이용을 위하여 필요한 부대시설 및 편익시설은 기반시설에 해당한다.
\end{choices}
\end{question}
\answer{5}
